\subsection{From sampled behavior to the GRPO objective}

Rollouts are generated by a frozen behavior snapshot $\pi_{\mathrm{old}}$, while gradients
update $\pi_\theta$. Importance sampling exactly relates the desired new-policy expectation
to old-policy trajectories through the response ratio
$\pi_\theta(o_i\mid q)/\pi_{\mathrm{old}}(o_i\mid q)$. For a long autoregressive response,
that exact ratio is a product of token ratios and can become extremely small or large.
GRPO therefore uses a local surrogate: each sampled token receives the trajectory
advantage and its own ratio
\begin{equation}
  \rho_{i,t}(\theta)
  =\frac{\pi_\theta(o_{i,t}\mid q,o_{i,<t})}
         {\pi_{\mathrm{old}}(o_{i,t}\mid q,o_{i,<t})}.
  \label{eq:compact-grpo-ratio}
\end{equation}
This per-token replacement is not algebraically equal to exact trajectory importance
sampling. It is a lower-variance local policy-improvement objective whose gradient at
$\pi_\theta=\pi_{\mathrm{old}}$ is the policy gradient in
Equation~\ref{eq:compact-policy-gradient}.

If a rollout batch is reused after an optimizer step, its actions become stale relative to
the changed policy. Clipping stops one old sample from justifying unlimited movement. A
separate frozen reference policy protects behavior outside the narrow reward, because many
small batch-relative updates can still accumulate into large drift. With
$d_{i,t}=\log\pi_{\mathrm{ref}}(o_{i,t})-\log\pi_\theta(o_{i,t})$, the complete loss is
\begin{equation}
  \mathcal{L}_{\mathrm{GRPO}}
  =-\frac{1}{G}\sum_i\frac{1}{T_i}\sum_t
  \left[
    \min\!\left(\rho_{i,t}A_i,
      \operatorname{clip}(\rho_{i,t},1-\epsilon,1+\epsilon)A_i\right)
    -\beta\left(e^{d_{i,t}}-d_{i,t}-1\right)
  \right].
  \label{eq:compact-grpo-loss}
\end{equation}
Read from left to right, the objective rewards above-average sampled decisions, corrects
for the policy that generated them, limits what stale evidence may justify, and penalizes
drift from the longer-lived reference model. The $1/T_i$ reduction gives each response
equal total normalization rather than allowing length alone to determine its weight.
